\documentclass[UTF8,11pt]{ctexart}
\usepackage{amsmath,amssymb,amsthm,mathtools}
\usepackage{booktabs,longtable,array,geometry,enumitem}
\usepackage{xcolor}
\definecolor{linkblue}{rgb}{0.10,0.20,0.45}
\usepackage[colorlinks=true,linkcolor=linkblue,citecolor=linkblue,urlcolor=linkblue]{hyperref}
\geometry{a4paper,margin=2.3cm}
\setlength{\parskip}{4pt}
\setlength{\emergencystretch}{3em}
\sloppy
\newcommand{\Lam}{\Lambda}
\newcommand{\cst}{c_\ast}
\newcommand{\ok}{\textcolor[rgb]{0.05,0.45,0.15}{\bfseries [已核]}}
\newcommand{\warn}{\textcolor[rgb]{0.72,0.28,0.02}{\bfseries [须处理]}}
\newcommand{\gap}{\textcolor[rgb]{0.55,0.10,0.10}{\bfseries [缺口]}}

\title{\bfseries 第二次查新复核报告\\[4pt]
\large 前三大特征值之和的精确上界\;$\Lam_3=\lambda_1+\lambda_2+\lambda_3$\\[2pt]
\normalsize 兼及副线 $\lambda_1+2\lambda_2$ 的极图分类}
\author{Haoyang Ji}
\date{2026 年 11 月 3 日}

\begin{document}
\maketitle

\section*{摘要与结论}

本报告记录 2026 年 11 月 3 日对主论文
\emph{A joint moment bound for the three largest graph eigenvalues}
所做的一次独立查新复核，以及据此对论文正文实施的修改。三条结论如下。

\begin{enumerate}[label=(\arabic*),leftmargin=2.2em]
\item \ok\ \textbf{主定理的新颖性安全。}
在可检索的主流文献（arXiv math.CO、OpenAlex 索引的期刊与学位论文、Zenodo）中，
\emph{没有任何工作改进过} Sivashankar--Tang--Wakhare（下称 STW，arXiv:2608.02429）
Proposition 6.1 给出的系数 $\frac{3+\sqrt5}{4}=1.309016994374947\ldots$。
本文的
\[
 c=\frac{106106902139519}{107310156800000}+\frac{41\sqrt3}{240}-\frac{53}{9000000}
  =1.284673255723918505\ldots
\]
仍是公开可查的最好系数，且证明是无条件的、自足的。

\item \warn\ \textbf{但下界常数已被一条时间更早、带 DOI 的自存档记录提出，必须引用。}
Zenodo 上署名 ``Tony Newton / Newton Astro Labs'' 的两条记录
（doi:10.5281/zenodo.21785906，2026-08-04；doi:10.5281/zenodo.22082842，2026-08-24）
以 ``pole--pentagon law''（极点--五边形律）的名义给出了与本文 \emph{完全相同}的常数与权重，
并进一步声称 $\Lam_3(G)\le\alpha_3n-3$。
该声称\textbf{依其自身表述是条件性的}：定理标题为 \emph{Candidate} universal $S_3$ theorem，
依赖一条机器证书链，作者明确要求独立审计。
论文已新增注 2.3 与两条参考文献如实处理，\textbf{本文不主张该常数的优先权}。

\item \ok\ \textbf{副线 $\lambda_1+2\lambda_2$ 未被抢先。}
所有针对加权谱和的标题级与全文级检索均无命中；
最近邻的三篇（树的凸组合、$S_2$ 精确解、阈图谱综述）都不覆盖本问题。
该章节可安全作为学位论文的独立成果。
\end{enumerate}

\section{检索方法、覆盖范围与局限}

\textbf{数据源。} arXiv 官方 API（\texttt{export.arxiv.org}，摘要级与全文级检索式）；
OpenAlex（\texttt{api.openalex.org}，\texttt{title.search} 与 \texttt{search} 两种过滤）；
Zenodo API（记录元数据、DOI 解析、版本关系、下载与浏览统计）；
arXiv 原文 LaTeX 源码（对关键近邻逐句核对）。
Crossref 在本次会话中不可达，Semantic Scholar 间歇性失败，
故以 OpenAlex 为主索引、以 arXiv 原文为最终裁决。

\textbf{检索式（节选）。} 主线：\texttt{"third spectral sum"}、
\texttt{"sum of the largest"} + \texttt{math.CO}、
\texttt{lambda\_1 lambda\_2 extremal}、\texttt{weighted spectral sum graphs}、
\texttt{pole pentagon spectral sum}、\texttt{three block graph eigenvalues}。
副线：\texttt{lambda\_1+2\textbackslash lambda\_2}、
\texttt{convex combination first and second eigenvalues}、
\texttt{stair graph eigenvalues extremal}、\texttt{Pell equation spectral graph theory}、
\texttt{second largest eigenvalue maximum graphs exact}。

\subsection{覆盖率的直接证据}

除检索式本身之外，本次还用四条独立途径交叉验证``没有后续文献''这一判断。

\begin{enumerate}[label=(\alph*),leftmargin=2.2em]
\item \textbf{引用图谱为空。} STW 在 OpenAlex 中的记录为
\texttt{W7172451143}（\emph{Graph Eigenvalues and Projection Constants}），
其 \texttt{cited\_by\_count} 为 \textbf{0}；
按 \texttt{filter=cites:W7172451143} 反查亦返回 0 条。
Huang--Wei（2607.23081）、KLMPT（2604.00512）、Paulin（2609.37980）三篇的
\texttt{cited\_by\_count} 同样为 \textbf{0}。
也就是说，这一方向在索引层面\emph{尚无任何后续工作}，
``已被别人改进''只可能来自未被学术索引收录的自存档。

\item \textbf{其他预印本平台为空。} OSF Preprints 标题检索
（\texttt{third spectral sum}、\texttt{spectral sum graphs}、
\texttt{eigenvalue sum graph}）三项均返回 \textbf{0 条}。
HAL 以 \texttt{title\_t} 检索 \texttt{"third spectral sum"} 返回 0 条，
检索 \texttt{"spectral sum"} 返回 7 条，逐条核对后全部属 QCD 谱求和规则
或``Spectral Sum Estimation''的计算复杂度论文，与图论无关。
viXra 的检索端点本次返回 404，SSRN 连接超时；
但两者均被 OpenAlex 收录，而 OpenAlex 的标题级检索对同一批检索式返回 0 条。

\item \textbf{``pole--pentagon''这一命名只存在于那两条 Zenodo 记录。}
OpenAlex 全文检索 \texttt{"pole-pentagon"} 命中 \textbf{2 条}，
即 doi:10.5281/zenodo.21785906 与 22082842，别无其他；
arXiv 全文检索 \texttt{all:"pole-pentagon"} 命中 \textbf{0 条}。
这说明该术语未进入任何同行评审文献或主流预印本。

\item \textbf{投影常数方向的 arXiv 全集只有 5 篇。}
以 \texttt{abs:"projection constant" AND cat:math.CO} 检索，
返回 2609.03200、2608.02429、2603.29280、2502.13752、1901.07866 共 5 篇，
其中前三篇已列入表~\ref{tab:main}
（2603.29280 与 2608.02429 标题相同，为同一工作的前身版本），
后两篇分别是凸体极化常数与极大投影常数的数值计算，与图论求和无关。

\item \textbf{新识别到的一篇期刊论文不影响结论。}
OpenAlex 检索中另出现
\emph{Extremal graphs for the sum of two largest eigenvalues}，
\textbf{AIMS Mathematics}（2026）。该文处理的是\emph{两项}和
$\lambda_1+\lambda_2$ 的极图，属 Huang--Wei 所在的同一问题族，
既不涉及 $\Lam_3$，也不涉及权重为 2 的 $\lambda_1+2\lambda_2$。
因 OpenAlex 在本次会话后段触发 IP 级限流（HTTP 429），
其 DOI 与作者名单暂未取到，留作后续补录。
\end{enumerate}

\textbf{局限（须向导师说明）。}
(i)~Zenodo 的站内关键词检索接口在本次会话中已失效
（任意查询均返回全站四万余条记录），因此 Zenodo 一侧的排查改为
``OpenAlex 标题索引反查 + 已知 DOI 直接核验''的组合，
不能排除 Zenodo 上存在未被 OpenAlex 收录的同类自存档。
(ii)~尚未出版的会议论文与在投期刊稿件不可检索。
(iii)~中文文献库（CNKI、万方）未纳入本次检索。

\section{主线：$\Lam_3$ 上界的文献比对}

表~\ref{tab:main} 汇总了所有与 $\Lam_3$ 上界相关、且本次已核验原文或摘要的工作。
其中 Tang（2603.21181）、Sivashankar（2603.28738）、Wakhare（2603.29280）
三项早期预印本的结果均被 STW 的 Proposition 6.1 涵盖；
本文的 $c$ 严格优于 STW。

\begin{table}[htbp]
\centering\small
\caption{与 $\Lam_3$ 上界相关的已核验文献}\label{tab:main}
\begin{longtable}{@{}p{4.6cm}p{6.4cm}p{4.2cm}@{}}
\toprule
文献 & 结果 & 与本文关系\\
\midrule
\endfirsthead
\toprule
文献 & 结果 & 与本文关系\\
\midrule
\endhead
Ebrahimi, SFU 硕士论文 (2008), Thm 3.4.1 &
$\tau_3\le3/\sqrt5=1.341640786\ldots$（补图 + Weyl + Cauchy--Schwarz） &
本文已引，为系数链第一环\\
Tang, arXiv:2603.21181 & $\lambda_3\le n/3-1$ & 单项，被 STW 涵盖\\
Sivashankar, arXiv:2603.28738 & $\Lam_3$ 型上界 & 被 STW 涵盖\\
Wakhare, arXiv:2603.29280 & $\Lam_3$ 型上界 & 被 STW 涵盖\\
KLMPT, arXiv:2604.00512 & $\lambda_1+\lambda_2\le8n/7$ & 两项和的普适上界\\
STW, arXiv:2608.02429, Prop.\ 6.1 &
$\Lam_k\le\frac{1+\gamma(k)}2n-k$；$k=3$ 给出 $\frac{3+\sqrt5}4n-3$；
并引入混合投影常数 $\eta_{\mathrm{mix}}(k)$ &
\textbf{本文的直接改进对象}\\
\textbf{Huang--Wei, arXiv:2607.23081} &
$S_2=\lambda_1+\lambda_2$ 的\textbf{精确}最大值与\textbf{唯一}极图 $K_n^\star$
（三部图，部大小由 $n\bmod7$ 决定），$S_2\le8n/7-2$，$7\mid n$ 时取等；
证明用 Ky Fan 变分 + 加权 Ferrers 商的谱不等式 &
两项问题已解决；\textbf{本次已补引}\\
\textbf{Paulin, arXiv:2609.37980} &
$\lambda_1,\lambda_2$ 的锐二次不等式与 $\lambda_2,\lambda_3$ 的二次不等式；
推出 $\lambda_1+\lambda_2\le8n/7-2$、$\lambda_3\le n/3-1$、$\lambda_2+\lambda_3\le2n/3-2$；
确定两组归一化特征值对的闭包 &
最近邻的二次不等式路线；\textbf{本次已补引}\\
\emph{Extremal graphs for the sum of two largest eigenvalues}, AIMS Mathematics (2026) &
两项和 $\lambda_1+\lambda_2$ 的极图 &
与 Huang--Wei 同族；DOI 待补（限流）\\
Sun--Guo--Ge--Das, arXiv:2609.23707 &
$S_k$ 的 Hoffman 型结果：刻画 $S_k<2k$ 的连通图，$P_n$ 为唯一极小图 &
属``极小/下界''方向，与本文不冲突\\
Kumar--Mohar--Mojallal--Pragada, arXiv:2609.03200 &
极大相对投影常数的稳定性 & 本文已引 \texttt{KMMP}\\
Leonida--Li, LAA, doi:10.1016/j.laa.2026.03.031 &
线性型幂的最小依赖 & 与本文 \texttt{LiLiu} 同方向\\
Schmidt, SJSU 硕士论文, doi:10.31979/etd.kkej-tu27 &
Powers 关于第三大特征值的猜想 & 单项 $\lambda_3$ 方向\\
\bottomrule
\end{longtable}
\end{table}

\section{须处理的先行记录：Zenodo ``Tony Newton''}

\subsection{记录本身}

\begin{itemize}[leftmargin=1.4em]
\item \textbf{记录 A}：doi:10.5281/zenodo.21785906，上传时间
2026-08-04T06:22:21Z，24 页，标题
\emph{The Third Spectral Sum of Graphs: An Exact Universal Bound,
Pole--Pentagon Extremality, and Resolution of the $S_3$ Extremal Problem}。
\item \textbf{记录 B}：doi:10.5281/zenodo.22082842，上传时间
2026-08-24T14:43:32Z，10 页，标题
\emph{The Third Spectral Sum of Graphs: Exact Projection Equivalence,
Pole--Pentagon Geometry, and Certified Carrier Reduction}。
\item \textbf{作者信息}：``Newton, Tony''，署名单位 ``Newton Astro Labs''，
联系邮箱 \texttt{tony.newton79@gmail.com}，ORCID 0009-0009-9189-6431；
API 返回的 \texttt{affiliations} 字段为空，即\textbf{无机构署名}。
许可 CC-BY-4.0。
\item \textbf{传播情况}：记录 A 下载 0 次、浏览 20 次（独立 19）；
记录 B 下载 2 次、浏览 8 次。被引数 0。
\item \textbf{账号行为模式}：同一账号批量上传彼此完全无关的主题
（Navier--Stokes、Erd\H{o}s--Straus、Goldbach 型问题等），
属典型的非策展式自存档，而非某个研究组的常规预印本发布。
\end{itemize}

\subsection{逐项重合点（已用 60 位十进制精度核验）}

\begin{center}\small
\begin{tabular}{@{}lll@{}}
\toprule
量 & 该记录的写法 & 本文的写法\\
\midrule
极值常数 &
$\alpha_3=\dfrac12+\dfrac{\sqrt5}{18}+\dfrac{\sqrt{35}}9=1.281568196872168\ldots$ &
$\cst=\dfrac{9+\sqrt5+2\sqrt{35}}{18}$\\[2pt]
顶点权重 & $p^\ast=\dfrac{91-25\sqrt7}{126}$ & $a=1-5b$\\[2pt]
环点权重 & $r^\ast=\dfrac{7+5\sqrt7}{126}$ & $b=\dfrac{7+5\sqrt7}{126}$\\[2pt]
极图模板 & axial pole $+$ regular pentagonal orbit & $K_1\vee C_5$ 六部核\\
\bottomrule
\end{tabular}
\end{center}

两个常数在代数上恒等：
$\frac12+\frac{\sqrt5}{18}+\frac{\sqrt{35}}9
 =\frac{9+\sqrt5+2\sqrt{35}}{18}$，
数值同为 $1.28156819687216787675243501399191724068\ldots$；
权重 $r^\ast$ 与本文的 $b$ 逐字相同，且 $p^\ast+5r^\ast=1$ 与本文 $a+5b=1$ 一致。
``极点--五边形''即本文的锥点 $K_1$ 加五边形环 $C_5$。

\subsection{该记录额外声称了什么，以及为什么它是条件性的}

记录 A 断言 $M_3(n)=\max_{|V(G)|=n}\Lam_3(G)\le\alpha_3n-3$ 对每个 $n$ 成立，
即它把 $\alpha_3$ 同时当作上界与下界，从而宣称 $S_3$ 极值问题已解决。
但同一份文本在三处明确限定该断言的地位：

\begin{quote}\small
\begin{enumerate}[label=(\alph*),leftmargin=1.8em]
\item 定理标题本身即为
``\emph{Theorem 1.1 (Candidate universal $S_3$ theorem -- exact certificate closure)}''；
\item 定理前提写作 ``\emph{Within the exact analytic and proof-carrying certificate chain}''；
\item 结论段落写作 ``\emph{Externally the responsible description remains
`Candidate Universal $S_3$ Theorem -- Exact Certificate Closure' until independent
specialists have audited the admissibility of the archived payments and replayed the
exact certificate stack}''，并在摘要中说明验收
``\emph{properly remains subject to independent mathematical audit of the admissibility
of the archived dual payments, the global sector partition and the exact verifier
implementation}''。
\end{enumerate}
\end{quote}

记录 B 同样把证书库描述为 ``\emph{when assumed complete}''，
并自设一节 ``\emph{Independent audit target}''。
换言之，\textbf{作者本人不主张该上界已被无条件证明}。
此外，两条记录都不含 graphon 方法、不含``每个极大化 graphon 的典范律至多 24 个支撑点''
这类结构定理，也不含本文的 $c=1.284673255723918505\ldots$。

\subsection{时间线与优先权判断}

本仓库 $\Lam_3$ 研究的首次提交为 \texttt{308390d}（2026-10-27），
晚于记录 A 约 \textbf{2.5 个月}、晚于记录 B 约 \textbf{2.1 个月}。
因此在下界常数 $\cst$、最优权重 $b$ 与 $K_1\vee C_5$ 模板这三项上，
\textbf{本文不构成首次提出}。而在无条件上界 $c$、
graphon 极大化的 24 支撑点结构定理、以及全部证书化正性验证这三项上，
未见任何在先公开材料，\textbf{本文构成首次}。

\subsection{已落实的处置}

\begin{enumerate}[label=(\roman*),leftmargin=2.2em]
\item 正文 §1 系数链之后新增两段：一段交代 Huang--Wei 与 Paulin 两项相邻结果
（说明为何三项和需要不同几何），一段直接点名 Newton 的两条 Zenodo 记录、
其条件性质、以及本文不主张常数优先权。
\item 引理 2.2（\texttt{attainment}）证明之后新增
\textbf{注 2.3}（\texttt{rem-newton}），标题即``$\cst$ 的优先权''，
逐字引用该记录对自身的限定语，并说明本文的下界只依赖 join 的常规商矩阵计算，
不需要其证书机制。
\item 第 10 节（Further questions）把无条件缺口
\[
 c-\cst=0.003105058851750628\ldots
\]
显式列出，并说明：若该记录的证书论证日后通过独立审计，则 $\mathcal M=\cst$，
区间收缩为一点；本文的定理与该结果无关。
\item 参考文献新增 4 条：\texttt{HW}、\texttt{Newton1}、\texttt{Newton2}、\texttt{Paulin}；
两条 Newton 条目在文献项内即注明``Zenodo 自存档、未经同行评审、结论条件性''，
使任何读者在参考文献表中就能看到限定。
\item 论文日期更新为 November 3, 2026；篇幅由 20 页增至 22 页。
\end{enumerate}

\subsection{风险评级}

\begin{center}\small
\begin{tabular}{@{}p{4.2cm}p{2.0cm}p{8.6cm}@{}}
\toprule
风险 & 等级 & 说明\\
\midrule
主定理被抢先 & \textcolor[rgb]{0.05,0.45,0.15}{低} &
无任何无条件改进 STW 的公开结果；本文的 $c$ 仍是公开最好系数\\
下界常数被抢先 & \textcolor[rgb]{0.72,0.28,0.02}{已发生，已处置} &
通过引用 + 明确让渡优先权 + 强调本文下界证明独立，风险已消除\\
被指学术不端 & \textcolor[rgb]{0.05,0.45,0.15}{低} &
论文现在主动、显著地披露了这条在先记录，而非回避\\
$\alpha_3$ 最终被证实 & \textcolor[rgb]{0.05,0.25,0.55}{不影响本文} &
本文从未主张 $c$ 最优；届时结论变为``本文给出无条件上界，
真实值为 $\cst$，缺口 $0.003105\ldots$ 由他人闭合''\\
$\alpha_3$ 被证伪 & \textcolor[rgb]{0.05,0.45,0.15}{不影响本文} &
本文的下界与上界均不依赖它\\
\bottomrule
\end{tabular}
\end{center}

\section{副线：$\lambda_1+2\lambda_2$ 的极图分类}

\subsection{检索结果}

针对该问题设计的 10 组 OpenAlex 标题级检索全部返回 \textbf{0 条}：
\texttt{weighted spectral sum graphs}、\texttt{lambda1 plus two lambda2}、
\texttt{exact maximum weighted combination adjacency eigenvalues}、
\texttt{first and second largest eigenvalues graphs}、
\texttt{eigenvalue sum extremal graph exact}、\texttt{three block graph eigenvalues}、
\texttt{stair graph eigenvalues extremal}、
\texttt{second largest eigenvalue maximum graphs exact}、
\texttt{Pell equation spectral graph theory}、
\texttt{spectral extremal problem weighted eigenvalues}。
arXiv 侧以 \texttt{all:"lambda\_1+2\textbackslash lambda\_2"} 检索，
唯一命中是一篇一维随机 Landau--Lifshitz 链的物理论文，与图论无关；
以 \texttt{abs:"spectral sum"} + \texttt{math.CO} 检索仅得 4 篇，
全部已在表~\ref{tab:main} 中列出。
Zenodo 侧以 \texttt{pole pentagon spectral sum} 反查，仅得上述两条 Newton 记录，
无任何 $\lambda_1+2\lambda_2$ 条目。

\subsection{最近邻文献与其不覆盖的理由}

\begin{center}\small
\begin{longtable}{@{}p{4.4cm}p{5.4cm}p{5.0cm}@{}}
\toprule
文献 & 结果 & 为何不构成抢先\\
\midrule
\endfirsthead
\toprule
文献 & 结果 & 为何不构成抢先\\
\midrule
\endhead
Kumar--Mohar--Pragada, arXiv:2601.10036v2（v1 2026-01-15，v2 2026-08-07） &
对\textbf{树}确定 $\alpha\lambda_1+(1-\alpha)\lambda_2$ 的极大与极小极值树，
$\alpha\in[0,1]$ &
类限制为树；且 $\lambda_1+2\lambda_2=3\bigl(\tfrac13\lambda_1+\tfrac23\lambda_2\bigr)$
中 $\lambda_2$ 的权重为 2，超出其凸组合范围\\
Huang--Wei, arXiv:2607.23081 &
$S_2$ 的精确最大值与唯一极图 & $k=1$ 情形；方法学最接近，结论不同\\
Manna--Mandal--Tura, arXiv:2610.22838（2026-10-16） &
阈图（$\{P_4,C_4,2K_2\}$-free）在邻接、Laplacian、符号 Laplacian、距离、
距离 Laplacian、距离符号 Laplacian、偏心、Seidel、归一化等 7 类矩阵下的谱综述
+ 开放问题清单 &
综述，无 $\lambda_1+2\lambda_2$ 结果；但阶梯图即阈图，\textbf{须作为该类图的标准引用}\\
arXiv:1602.02069 & cograph 与 $P_5$-free 图的谱性质、特征值重数上界猜想 &
背景引用\\
Paulin, arXiv:2609.37980 & 归一化特征值对的区域闭包 & 区域型结果，未涉及权重 2\\
\bottomrule
\end{longtable}
\end{center}

\subsection{章节写作时的引用义务}

在学位论文的 $\lambda_1+2\lambda_2$ 章节中，以下四点必须写明：
(i)~把 $C=\frac{5+4\sqrt2}{7}=1.522407\ldots$ 与 Huang--Wei 的 $\frac87$ 并列，
说明权重 2 使系数跨过 1.5；
(ii)~指出 Kumar--Mohar--Pragada 只处理树且只处理凸组合，本问题是全体图上的非凸加权；
(iii)~把阶梯图类明确等同于阈图，并引用 arXiv:2610.22838 综述；
(iv)~说明 Pell 型并列 $k^2-2n^2=-7$ 是本问题特有的算术现象，
在已检索文献中未见对应先例。

\section{给导师的一句话说明}

本文对 $\Lam_3$ 给出的无条件上界 $cn-3$ 是公开的最好结果，未被抢先；
下界常数 $\cst$ 与 $K_1\vee C_5$ 极图模板在本文之前约两个半月
由一条 Zenodo 自存档记录以``极点--五边形律''的名义提出过，
该记录同时声称（条件性地、依赖未经审计的机器证书）这就是精确值，
本文已在注 2.3 与参考文献中如实披露并让渡该常数的优先权，
主定理的证明不依赖它；副线 $\lambda_1+2\lambda_2$ 经检索未被抢先。

\section{本次查新的原始数据}

全部抓取结果与检索脚本已入库，路径
\texttt{paper/literature/novelty\_20261103/}：
\texttt{raw\_hits.json}、\texttt{fast\_hits.json}、\texttt{lam12lam2\_hits.json}、
\texttt{round3\_log.txt}、\texttt{round4\_log.txt}，
以及 \texttt{\_search\_lam12.py}、\texttt{\_search\_round2.py}、
\texttt{\_search\_round3.py}、\texttt{\_search\_round4.py}。
两条 Newton 记录的 PDF 与纯文本抽取为
\texttt{paper/literature/Newton2026a.pdf}/\texttt{.txt} 与
\texttt{Newton2026b.pdf}/\texttt{.txt}；
STW 与 Huang--Wei 的 arXiv LaTeX 全文分别为
\texttt{STW\_src.tex.txt} 与 \texttt{HuangWei\_src.tex.txt}。

\end{document}
